%% tables/tbl_class_representatives.tex
%% Table II — 분류군별 대표 기법 + 메커니즘 (table* + resizebox, multirow 제거)

\begin{table*}[t]
\centering
\caption{5 분류군 대표 기법의 단독 효과와 실패/성공 메커니즘. 각 분류군은 자원의 도메인에 따라 정의되며, 동일 분류군 내 기법들은 공통 메커니즘으로 실패 또는 성공한다.}
\label{tbl:class-representatives}
\scriptsize
\setlength{\tabcolsep}{4pt}
\resizebox{\textwidth}{!}{%
\begin{tabular}{@{}cp{4.3cm}rlp{6.0cm}@{}}
\toprule
\textbf{Class} & \textbf{Representative} & \textbf{$\Delta$ (1-run)} & \textbf{Verdict} & \textbf{Mechanism} \\
\midrule
A & AMX matmul drop-in              & $+0.26\%$  & noise         & H2D/D2H sync overhead 누적이 microbench gain 상쇄 \\
A & 4-method Jacobi router          & $-94.20\%$ & catastrophic  & LM-head 1.81\,s $\times$ 165 chat req 직렬 큐 collapse \\
B & NUMA pin (dual instance)        & $\mathbf{+1.54\%}$ ($\uparrow$ multi $+2.24\%$) & \textbf{robust} & local/remote 2.1$\times$ latency 비대칭의 hardware-level 회수 \\
B & 4-OS-lever stack                & $-0.03\%$  & destructive   & 동일 도메인 stacking $\rightarrow$ context-switch + cache thrash \\
C & task-pool fill (actual fire)    & $-1.75\%$  & catastrophic  & step-internal lock contention 으로 critical path stall \\
D & AMX draft proxy (Qwen 0.5B)     & $-2.77\%$  & failure       & microbench 0.524\,ms, invasive integration 부재로 e2e 변환 실패 \\
E & branchy $\times$ 100\,ms (cellB) & $\mathbf{+5.28\%}$ (cell $+12.04\%$) & \textbf{threshold} & step-boundary 정렬 + prefetcher 억제의 결합 효과 \\
E & regular $\times$ 10\,ms (cellA)  & $+0.98\%$  & noise         & 비정수배 alignment 로 critical-path 침범 \\
\bottomrule
\end{tabular}%
}
\smallskip
\par\noindent\scriptsize 본 분류군 대표 검증은 가설 H1--H3 (\S\ref{subsec:hypotheses}) 의 정당화 근거. 분류군 B 와 E 의 cross-domain pair (NUMA $\oplus$ branchy) 가 \algname 의 \textsc{Chord} + \textsc{Accent} 단계 design 의 직접 backing (Table~\ref{tbl:superposition}).
\end{table*}
